\documentclass[10pt]{article}
\usepackage{vendor-kestrel}
\renewcommand{\DocID}{BOARD-001}
\renewcommand{\RunningPart}{BLENDER-8}
\hypersetup{pdftitle={Kestrel Blender-8 / Control Chassis Integration and Operating Requirements}}
\begin{document}
\VendorTitle{BLENDER-8}{Control chassis integration}{Command assembly / controller / display / motor drive / case-temperature acquisition}
\MetaStrip{SYSTEM INTERFACE 02\hfill ISSUE P2\hfill 06 OCT 2026}
\section*{Equipment configuration}
The control chassis combines a Northstar B8 controller, BA-8 command assembly, TriAxis MX8-1 input selector, PixelRiver PX32-16 display, Vortek MD20 drive and ThermaSense AVT10 case sensor.
\begin{center}
\begin{tikzpicture}[font=\scriptsize,>=Stealth]
\node[draw,fill=VendorTint,minimum width=27mm,minimum height=12mm,align=center](m)at(0,0){B8\\controller};
\node[draw,minimum width=25mm,minimum height=12mm,align=center](mx)at(-4.5,0){MX8-1\\selector};
\node[draw,minimum width=25mm,minimum height=12mm,align=center](ba)at(-4.5,2){BA-8\\commands};
\node[draw,minimum width=27mm,minimum height=12mm,align=center](lcd)at(4.5,2){PX32-16\\display};
\node[draw,minimum width=27mm,minimum height=12mm,align=center](motor)at(4.5,-1){MD20\\motor/driver};
\node[draw,minimum width=27mm,minimum height=12mm,align=center](temp)at(0,-2.7){AVT10\\case sensor};
\draw[->](ba)--node[left]{C0--C7}(mx);
\draw[->](mx.east)--node[above]{Y / PA3}(m.west);
\draw[->](m.south west)--++(-.3,-.3)-|node[pos=.42,below]{PA0--PA2 / S0--S2}(mx.south);
\draw[<->](m.north east)--node[above,sloped]{XBUS / VBLANK}(lcd.south west);
\draw[->](m.east)--node[above,sloped]{PWM / enable gate}(motor.west);
\draw[->](motor.north)--(4.5,1.0)--(0,1.0)--(m.north);
\node[above]at(.6,1.02){TACH / PB1};
\draw[->](ba.east)--(-1.4,2.0)--(-1.4,.65)--(m.north west);
\node[above]at(-2.3,2.0){STOP\_N / PB2};
\draw[->](temp)--node[right]{VOUT / AN0}(m);
\draw[dashed](motor.south)|-node[pos=.7,below]{case attachment}(temp.east);
\end{tikzpicture}
\end{center}
\par\noindent\begin{tabularx}{\linewidth}{@{}l Y@{}}
\toprule\TableHead Board connection & Signal ownership\\\midrule
PA0, PA1, PA2 & Controller outputs to S0, S1, S2; S0 is least significant.\\
PA3 & Input from MX8-1 Y; never configure as opposing output.\\
C0--C6 / C7 & Speed 1--7 / PULSE to X0--X7; board pull-ups; active low.\\
PB2 & Dedicated active-low STOP\_N, separate from mux.\\
PB0 & Firmware-owned MOTOR\_EN; board pull-down at reset.\\
Dedicated PWM & Output to MD20 PWM.\\
PB1 & MD20 tach output to hardware edge counter.\\
AN0 / AN1 & AVT10 VOUT with diagnostic pull-up / grounded spare.\\
XBUS / VBLANK & Display register transfers / dedicated blanking-edge input.\\\bottomrule
\end{tabularx}
\begin{Notice}{ACTUAL DRIVE ENABLE}
MD20 ENABLE = panel power AND BA-8 RUN\_PERMIT AND PB0. STOP removes the independent permission regardless of firmware state. Verify the gate separately from the firmware-owned PB0 and PWM outputs.
\end{Notice}
\vfill\Caption{Interconnect is functional. Layout, connector pin numbering and physical power-distribution details are not defined by this drawing.}
\clearpage
\Continuation{Operating and input requirements}{Controlled behavior at startup, selection, PULSE and STOP.}
\section*{Power-on qualification}
Initialize MOTOR\_EN, PWM enable and pending duty to zero before enabling interrupts or indicating readiness. Remain disabled until fresh, plausible case-temperature samples stay at or below 65 $^\circ$C for 2 s, STOP is released, and all eight command contacts have been released for at least 20 ms.

A speed retained across power loss cannot initiate motion. Only a later qualified command can start the drive. Mechanical latches, shaft momentum and retained heat are not cleared by an electrical reset.
\section*{Input acquisition}
Allow at least 2 $\mu$s for mux settling after final address selection. Sample every command contact at least once each 8 ms. Accept transitions only after consistent observations span at least 16 ms. Stable commands must be recognized within 40 ms. PULSE holds shorter than 50 ms need not be recognized.

Two or more speed contacts persisting for at least 40 ms are an invalid-input fault. Short scan-skew or bounce transients must not select an arbitrary winning speed or create an immediate persistent fault.
\section*{Nominal speed requests}
\begin{center}
\begin{tabular}{@{}c r c r@{}}
\toprule\TableHead Selection & Target RPM & Selection & Target RPM\\\midrule
1 & 3,000 & 5 & 13,000\\
2 & 5,500 & 6 & 15,500\\
3 & 8,000 & 7 & 18,000\\
4 & 10,500 & Idle & Drive off\\\bottomrule
\end{tabular}
\end{center}
Use the MD20 nominal unloaded characteristic to establish the requested drive fraction. At nominal no load, one second after selection, speed error must not exceed the greater of 100 RPM or 2\% of target. Loaded speed sag is allowed; closed-loop regulation is not a system requirement.

PULSE adds 2,000 RPM to the selected request, capped at 20,000 RPM. From idle it requests 2,000 RPM. Release restores the selected speed or idle. PULSE is neither 10\% PWM nor a periodic chopping mode.
\section*{Stop and indication}
STOP independently inhibits the drive and releases commands. Firmware must also clear MOTOR\_EN, PWM enable and pending duty within 10 ms of STOP\_N assertion. Release never restores a prior run request.

The pixel display shows 0 for idle, 1--7 for speed, P for standalone PULSE, or 1P--7P for boosted speed. Generate the patterns through the display's pixel RAM. Present the dominant latched fault within 100 ms, measured at the scanned display, without delaying drive removal.
\vfill\Caption{Control response and shutdown are measured at commanded outputs, not at the instant the rotor reaches zero RPM.}
\clearpage
\Continuation{Supervisory limits and recovery}{Independent latched causes / deliberate acknowledgment / no automatic restart.}
\section*{Acquisition and fault qualification}
Evaluate protection at least every 10 ms. Take AN0 conversions at intervals no greater than 10 ms and qualify freshness from completed conversions. For B8 nominal reference, use $V=3.3N/1023$ and $T=(V-0.500)/0.010$ in signed temperature arithmetic.
\par\noindent\begin{tabularx}{\linewidth}{@{}l l Y@{}}
\toprule\TableHead Label & Mask & Qualifying condition\\\midrule
TS & 04h & A fresh AN0 code below 62 or above 574; or no completed sample for more than 100 ms.\\
TH & 02h & Two consecutive fresh, plausible readings at or above 85 $^\circ$C.\\
IF & 08h & Two or more speed contacts persist for at least 40 ms.\\
ST & 01h & Unverified motion meets the grace and persistence conditions below.\\\bottomrule
\end{tabularx}
Remove commanded drive within 10 ms after fault qualification. Retain every independently qualified cause. Display priority is TS, then TH, then IF, then ST. An implausible rail code is not also trustworthy evidence for TH. TH means full thermal shutdown, not reduced drive.
\section*{Unverified-motion supervision}
Estimate speed from tach-count differences over a rolling 200 ms interval, using two rising edges per revolution and modulo-65,536 count differences. Track the last observed count change.

Allow a 500 ms grace interval only on an actual commanded off-to-on transition. Speed changes and an added PULSE while already running do not restart it. Thereafter, while drive is commanded on, a candidate exists when estimated speed is below 300 RPM OR no count change has been observed for at least 150 ms. A candidate persisting for 200 ms latches ST.

A startup with no shaft motion must be disabled within 750 ms of command. A sudden lock after established running must be disabled within 400 ms. Thermal and signal-validity protection remain active during startup grace and after shutdown. Intentional coastdown with drive off does not create a new ST.
\section*{Recovery sequence}
While faulted, require fresh, plausible readings at or below 65 $^\circ$C continuously for 2 s and no current invalid input combination. Then require a \textbf{new} STOP press, stable for 20 ms. STOP must subsequently be released, with every command contact released for 20 ms, before READY. Only a later new command starts motion.
\begin{Notice}{NO QUEUED ACKNOWLEDGMENT}
Holding STOP while the unit cools does not pre-authorize recovery. Cooling alone, restored tach pulses, repaired sensor wiring and removal of an obstruction do not clear the fault latch. Automatic retries and reverse kicks are prohibited.
\end{Notice}
\vfill\Caption{Kestrel Appliance Systems / BOARD-001. Case sensing does not establish winding-hotspot temperature; ST does not uniquely identify a mechanical obstruction.}
\end{document}
